\documentclass[10pt,letterpaper,twocolumn]{article}
\usepackage[margin=0.65in]{geometry}
\usepackage[T1]{fontenc}
\usepackage[utf8]{inputenc}
\usepackage{lmodern}
\usepackage{enumitem}
\usepackage{xcolor}
\usepackage{booktabs}
\usepackage{array}
\usepackage{hyperref}
\setlist[itemize]{leftmargin=1.1em, itemsep=0.2em, topsep=0.2em}
\setlist[enumerate]{leftmargin=1.2em, itemsep=0.2em, topsep=0.2em}
\setlength{\columnsep}{0.25in}
\setlength{\parindent}{0pt}
\setlength{\parskip}{0.25em}

\newcommand{\sectiontitle}[1]{\vspace{0.35em}\textbf{\large #1}\vspace{0.15em}}
\newcommand{\subhead}[1]{\textbf{#1}}

\title{\vspace{-1.2em}PR2 Team Lead Brief: Execution, Ownership, and Deliverables}
\author{COMP 4710 Group 11}
\date{March 5, 2026}

\begin{document}
\maketitle
\thispagestyle{empty}

\begin{center}
\small
\textbf{How to use in meeting:} move section by section, lock decisions live, and end only after each owner has a concrete output contract.\\
\textbf{Intent:} detailed ownership + methods + deliverables, minimal scheduling.
\end{center}

\sectiontitle{1) Non-Negotiable PR2 Scope}
\begin{itemize}
  \item Include all core analysis tracks: comparison, network, coverage, capacity, equity, walkability, and multi-city comparison.
  \item Keep report and Streamlit synchronized: same metrics, same claims, same figure logic.
  \item Final outputs must be regenerable through notebooks and reporting pipeline.
  \item README must reflect actual pipeline with Mermaid diagram (no legacy/incorrect flow diagrams).
  \item All claims in text must link to at least one figure/table.
\end{itemize}

\subhead{Out of scope unless critical}
\begin{itemize}
  \item New exploratory analysis that does not improve final recommendations.
  \item One-off figure generation paths outside notebook/reporting workflow.
\end{itemize}

\sectiontitle{2) Team Contracts (Owner by Owner)}

\subhead{Ahmed (Team Lead): Equity + Walkability + Final Synthesis}

\subhead{Mission}
\begin{itemize}
  \item Produce the final recommendation narrative: where inequity is highest, where walkability limits access, and what should be prioritized first.
\end{itemize}

\subhead{Methods/math}
\begin{itemize}
  \item Need index from vulnerability indicators (normalized + weighted).
  \item Service index from coverage/frequency/access indicators.
  \item Equity gap: \texttt{need\_index - service\_index}.
  \item Walkability index from normalized built-environment proxies.
  \item Final priority score combining equity gap, walkability deficit, and strategic importance.
\end{itemize}

\subhead{Libraries}
\begin{itemize}
  \item \texttt{pandas}, \texttt{numpy}, \texttt{geopandas}, \texttt{duckdb}, \texttt{matplotlib}.
\end{itemize}

\subhead{Deliverables}
\begin{itemize}
  \item Equity outputs by neighborhood/corridor (table + visuals).
  \item Walkability outputs tied to transit service.
  \item Combined priority list for interventions.
  \item Final policy implications narrative (report + presentation speaking points).
\end{itemize}

\subhead{Handoffs received}
\begin{itemize}
  \item Coverage priority outputs from Sudipta.
  \item Capacity/reliability priorities from Cathy.
\end{itemize}

\subhead{Handoffs produced}
\begin{itemize}
  \item Final prioritized recommendation matrix for full team packaging.
\end{itemize}

\subhead{Acceptance check}
\begin{itemize}
  \item Recommendations are evidence-backed, ranked, and traceable to figures.
\end{itemize}

\vspace{0.4em}
\subhead{Cathy: Capacity + Network Operations}

\subhead{Mission}
\begin{itemize}
  \item Quantify operational stress and reliability risk, then map each issue to an intervention type.
\end{itemize}

\subhead{Methods/math}
\begin{itemize}
  \item Capacity proxy: \texttt{pphpd = departures\_per\_hour * vehicle\_capacity * load\_factor}.
  \item Stress score from crowding, passups, delays, and reliability indicators.
  \item Priority score: stress + equity modifier + strategic corridor modifier.
  \item Tier diagnostics: compare stress/reliability across route tiers.
\end{itemize}

\subhead{Libraries}
\begin{itemize}
  \item \texttt{pandas}, \texttt{numpy}, \texttt{duckdb}, \texttt{matplotlib}, \texttt{networkx}, optional \texttt{seaborn}.
\end{itemize}

\subhead{Deliverables}
\begin{itemize}
  \item Route/corridor stress ranking outputs.
  \item Reliability summary outputs (tier/time perspective).
  \item Upgrade/intervention shortlist with rationale.
  \item Transfer/operational network insight summary.
\end{itemize}

\subhead{Handoffs received}
\begin{itemize}
  \item Final equity weighting guidance from Ahmed.
\end{itemize}

\subhead{Handoffs produced}
\begin{itemize}
  \item Capacity intervention table for synthesis + dashboard + report.
\end{itemize}

\subhead{Acceptance check}
\begin{itemize}
  \item Every intervention has a measurable trigger and supporting evidence.
\end{itemize}

\sectiontitle{3) Team Contracts (continued)}

\subhead{Sudipta: Coverage + Accessibility}

\subhead{Mission}
\begin{itemize}
  \item Identify underserved areas and produce actionable coverage improvement priorities.
\end{itemize}

\subhead{Methods/math}
\begin{itemize}
  \item Spatial joins of service geometry to neighborhood units.
  \item Coverage/accessibility thresholds based on service proximity + frequency.
  \item Optional segmentation/clustering for neighborhood transit profile groups.
  \item Underserved ranking combining low service and contextual need.
\end{itemize}

\subhead{Libraries}
\begin{itemize}
  \item \texttt{geopandas}, \texttt{pandas}, \texttt{numpy}, \texttt{duckdb}, \texttt{matplotlib}, optional \texttt{scikit-learn}.
\end{itemize}

\subhead{Deliverables}
\begin{itemize}
  \item Coverage outputs by neighborhood.
  \item Accessibility and underserved-area outputs.
  \item Coverage improvement priority outputs for report and app.
  \item Canonical neighborhood coverage table for reuse by all members.
\end{itemize}

\subhead{Handoffs received}
\begin{itemize}
  \item Shared metric definitions (need/service) with Ahmed.
\end{itemize}

\subhead{Handoffs produced}
\begin{itemize}
  \item Neighborhood coverage priority data package for synthesis and visuals.
\end{itemize}

\subhead{Acceptance check}
\begin{itemize}
  \item Coverage rankings are reproducible and align with source geographic data.
\end{itemize}

\vspace{0.4em}
\subhead{Stephenie: Visualization + Streamlit + Packaging}

\subhead{Mission}
\begin{itemize}
  \item Deliver one coherent visual layer so report, dashboard, and deck present identical conclusions.
\end{itemize}

\subhead{Methods/implementation}
\begin{itemize}
  \item One consistent map/figure standard (legend, scale, labeling, style).
  \item One figure registry with owner, section, filename, and status.
  \item Dashboard sections built from finalized metric outputs (no separate logic).
  \item Final consistency pass across all visual artifacts.
\end{itemize}

\subhead{Libraries}
\begin{itemize}
  \item \texttt{streamlit}, \texttt{matplotlib}, optional \texttt{plotly}, shared visualization helpers.
\end{itemize}

\subhead{Deliverables}
\begin{itemize}
  \item Final standardized figure package.
  \item Streamlit pages for all PR2 sections.
  \item Presentation visual kit.
  \item Updated README Mermaid pipeline diagram.
\end{itemize}

\subhead{Acceptance check}
\begin{itemize}
  \item No metric drift: dashboard and report match exactly.
\end{itemize}

\sectiontitle{4) Cross-Team Rules You Enforce as Lead}
\begin{itemize}
  \item \textbf{Single metric dictionary}: metric name, formula, unit, source table, interpretation.
  \item \textbf{Single figure registry}: owner, file, section, status, regeneration source.
  \item \textbf{Single evidence map}: every claim maps to figure/table reference.
  \item \textbf{Reproducibility gate}: output is not ``done'' until regenerable.
  \item \textbf{Consistency gate}: section definitions cannot conflict.
  \item \textbf{Narrative gate}: recommendations must be ranked and evidence-backed.
\end{itemize}

\subhead{Fast escalation protocol}
\begin{itemize}
  \item If a blocker appears, owner posts blocker + required input + who can unblock.
  \item If definitions conflict, lead resolves one final definition and freezes it.
  \item If integration quality degrades, freeze scope and finish consistency first.
\end{itemize}

\sectiontitle{5) Meeting Agenda (Detailed, but not schedule-heavy)}
\begin{enumerate}
  \item \textbf{Scope lock}: confirm non-negotiables and final outputs.
  \item \textbf{Ahmed contract}: confirm equity/walkability formulas, outputs, and final synthesis responsibilities.
  \item \textbf{Cathy contract}: confirm capacity stress logic, reliability interpretation, and intervention outputs.
  \item \textbf{Sudipta contract}: confirm underserved definitions, coverage outputs, and handoff format.
  \item \textbf{Stephenie contract}: confirm visual standards, Streamlit structure, and packaging conventions.
  \item \textbf{Cross-team governance}: lock dictionary/registry/evidence-map ownership.
  \item \textbf{Final lock}: each owner states final deliverables and acceptance criteria.
\end{enumerate}

\subhead{Questions to force clear commitments}
\begin{itemize}
  \item What is your exact output artifact list?
  \item What formula/logic drives each ranked result?
  \item What dependencies do you require from others?
  \item What handoff format are you giving and receiving?
  \item What is your acceptance test for ``done''?
\end{itemize}

\sectiontitle{6) Definition of Done (PR2)}
\begin{itemize}
  \item All required sections are complete and consistent.
  \item Capacity and walkability are both explicit in findings and recommendations.
  \item Multi-city comparison is included and interpretable.
  \item Streamlit, report, and deck use one consistent metric story.
  \item Final outputs regenerate from the agreed workflow.
\end{itemize}

\sectiontitle{7) Quick Close Script (Team Lead)}
\begin{itemize}
  \item ``Each owner is locked to specific outputs and acceptance tests.''
  \item ``If it is not reproducible and traceable, it is not done.''
  \item ``No conflicting metric definitions after this meeting.''
  \item ``Recommendations must be ranked and evidence-backed.''
\end{itemize}

\end{document}
